\documentclass[10pt,a4paper]{amsart}
\usepackage[margin=19mm]{geometry}
\usepackage{amsmath,amssymb,mathtools}
\usepackage[scr=rsfs,cal=euler]{mathalfa}
\usepackage{xfrac}
\usepackage[unicode,hidelinks,bookmarks=false]{hyperref}
\newcommand{\ie}{i.e.\ }
\newcommand{\cf}{cf.\ }
\newcommand{\resp}{respectively\ }
\newcommand{\Subsection}{\subsection}
\providecommand{\nobreakdash}{\nobreak\hbox{-}}
\setlength{\emergencystretch}{2em}
\multlinegap0pt
\numberwithin{equation}{section}
\usepackage{xcolor}
\newcommand*\dst{\displaystyle}
\newcommand*\eps{\varepsilon}
\newcommand*\p{\partial}
\renewcommand*\a{\alpha}
\newcommand*\R{\mathbb{R}}
\newcommand*\C{\mathbb{C}}
\newcommand*\M{\mathcal{M}}
\newcommand*\N{\mathbb{N}}
\newcommand*\f{\frac}
\newcommand{\clC}{\mathcal{C}}
\newcommand{\clB}{\mathcal{B}}
\newcommand{\comMD}[1]{\textcolor{black}{#1}}
\newcommand{\Mig}[1]{{\textcolor{black}{#1}}}
\newcommand{\Mag}[1]{{\textcolor{black}{#1}}}
\newcommand{\MT}[1]{{\textcolor{black}{#1}}}
\newcommand*\supp{{\mbox{\small\rm supp}}\,}
\newcounter{stepcount}
\newenvironment{step}[1]{\refstepcounter{stepcount}\begin{proof}[\textbf{Step~\thestepcount.\;{#1}}]}{\let\qed\relax\end{proof}}
\newcommand*{\mk}{\mkern -1mu}
\newcommand*{\Mk}{\mkern -2mu}
\newcommand*{\mK}{\mkern 1mu}
\newcommand*{\MK}{\mkern 2mu}
\newcommand*{\relabel}{\renewcommand{\labelenumi}{(\theenumi)}}
\newcommand*{\romanenumi}{\renewcommand*{\theenumi}{\roman{enumi}}\relabel}
\newcommand*{\Romanenumi}{\renewcommand*{\theenumi}{\Roman{enumi}}\relabel}
\newcommand*{\alphenumi}{\renewcommand*{\theenumi}{\alph{enumi}}\relabel}
\newcommand*{\Alphenumi}{\renewcommand*{\theenumi}{\Alph{enumi}}\relabel}
\newcommand*{\Hypenumi}{\renewcommand*{\theenumi}{\textbf{Hyp-\arabic{enumi}}}}
\let\oldtilde\tilde
\renewcommand*{\tilde}[1]{\mathchoice{\widetilde{#1}}{\widetilde{#1}}{\oldtilde{#1}}{\oldtilde{#1}}}
\let\oldexists\exists
\renewcommand*{\exists}{\mathrel{\oldexists}}
\let\oldforall\forall
\renewcommand*{\forall}{\mathrel{\oldforall}}
\newtheorem{defi}{Definition}[section]
\newtheorem{theo}[defi]{Theorem}
\newtheorem{prop}[defi]{Proposition}
\newtheorem{lemm}[defi]{Lemma}
\newtheorem{coro}[defi]{Corollary}
\newtheorem{rema}[defi]{Remark}
\title[Nonnegativity of the fragmentation power-series solution]{An inverse problem: recovering the fragmentation kernel from the short-time behaviour of the fragmentation equation: original Theorem 2.6}
\author{Marie Doumic, Miguel Escobedo and Magali Tournus}
\date{Source: Annales Henri Lebesgue 7 (2024), 621-671; DOI 10.5802/ahl.207}
\begin{document}
\maketitle
\noindent\textit{Editorial scope.} The counted unit is original Theorem 2.6 with its complete author proof. Hypotheses (Hyp-1) and (Hyp-2), equation (2.4), Definition 2.1, Theorem 2.2, Proposition 2.3, Theorem 2.4 with its full authored proof and Remark 2.5 are retained as uncounted context so that the statement and proof stand alone. Section 3 and the inverse-problem, stability, reconstruction and numerical material are omitted. Original theorem, equation and citation numbers are retained. Print slips, including the missing limit value in the BL convergence sentence, the lowercase $k_m$ written for $\kappa_m$ and the source two-word title 'Non negativity', are kept literal without generated corrections. Only the AMS wrapper supplies editorial formatting.
\setcounter{section}{1}
\section{Measure-valued solutions to the fragmentation equation: existence, uniqueness, stability and series representation}
\label{Section20}

The basis of our analysis in all the remaining of this work are the measure-valued solutions to the Cauchy problem for equation~\ref{eq:frag} with the initial condition
\begin{equation}\label{eq:data}
\mu _t(\MT{t=0})=\mu _0,
\end{equation}
whose precise definition is given below. Throughout the present paper, the following assumptions are used.
\begin{enumerate}\Hypenumi
\item \label{hyp1}
The fragmentation kernel $\kappa \in \M^+(\R^+)$ contains no atom at $x=0$ and at $x=1$, and satisfies
\begin{equation}
\supp(\kappa)\subset [0,1], \qquad \int\limits_0^1d\kappa(z) =N<+\infty, \qquad \int\limits_0^1 z d\kappa(z)=1.
\end{equation}
\item \label{hyp2}
The initial condition $\mu_0 \in \M^+(\R^+)$ is compactly supported
\begin{equation}
\supp(\mu_0)\subset [0,\MT{L}].
\end{equation}
\end{enumerate}


Even though $\kappa$ and $\mu_t$ are measures, we sometimes write the fragmentation equation as
\begin{equation}\label{eq:frag}
\begin{aligned}
\dfrac{\p}{\p t}\mu_t (x)& =- \a x^{\gamma}\mu_t(x) + \alpha \dst\int_{x}^{\infty}\kappa\left(\dfrac{x}{y} \right)y^{\gamma-1} d\mu_t(y),\qquad
\mu_{t=0}(x) = \mu_0(x),
\end{aligned}
\end{equation}
or as
\begin{equation*}\label{eq:frag2}
\begin{aligned}
\dfrac{\p}{\p t}\mu_t (x)& =- \a x^{\gamma}\mu_t(x) + \alpha \dst\int_{x}^{\infty}\kappa\left(\dfrac{x}{y} \right)y^{\gamma-1} \mu_t(y)dy,\qquad
\mu_{t=0}(x) = \mu_0(x).
\end{aligned}
\end{equation*}

\begin{defi}[Measure-valued solution for~\eqref{eq:frag}]\label{def:measure_valued_sol}
A family $(\mu_t)_{t\,\geq\,0} \subset \M\linebreak(\R^+)$ is called a measure-valued solution to problem~\href{https://ahl.centre-mersenne.org/item/AHL_2024__7__621_0/}{(1.3)}~\href{https://ahl.centre-mersenne.org/item/AHL_2024__7__621_0/}{(1.4)}~\eqref{eq:data} with initial data $\mu_0 \in \M(\R^+)$ satisfying~\eqref{hyp2} if the mapping $t\to \mu_t$ is narrowly continuous and for all $\varphi \in \clC(\R^+)$ such that $x\mapsto \varphi(x)/(1+x)$ is bounded on $[0,\infty)$, and all $t \geq 0$,
\begin{multline}\label{ppcm1}
\int_{\R^+} \varphi(x) d\mu_t(x)\\
= \int_{\R^+} \varphi(x) d\mu_0(x) + \dst\int_0^t ds \dst\int_{\R^+} d\mu_s(x) \alpha x^{\gamma}\left(- \varphi(x) + \dst\int_{0}^{1} d\kappa(z)\varphi(xz)\right).
\end{multline}
\end{defi}

We recall that $\mu_n$ converges narrowly toward $\mu$ if for all $\varphi \in C_b(\R^+)$, $\int \varphi d\mu_n \to \int \varphi d\mu$, where $C_b(\R^+)$ denotes the set of continuous and bounded functions defined on $\R^+$.

Although several results may be found in the references given in the introduction about the existence and uniqueness of solutions to fragmentation equations, none of them covers exactly the hypotheses that we have in mind for $\kappa$ and the initial data $\mu _0$. For the sake of completeness, our first result is then an existence and uniqueness of compactly supported and non negative measure-valued solutions to~\eqref{eq:frag} under assumptions~\eqref{hyp1}, \eqref{hyp2}. We begin with a uniqueness and stability result.

\begin{theo}[Uniqueness and TV-stability for the fragmentation equation in $(\M(\R^+),\|.\|_{TV})$]\label{thm:WP_BV}
Assume~\eqref{hyp1}, \eqref{hyp2} and $\comMD{\gamma\geq 0}$. Suppose that $\mu _t\in \clC(\R^+, \M(\R^+))$ is a measure-valued solution to~\eqref{eq:frag}, in the sense of Definition~\ref{def:measure_valued_sol}.
Then, for all $t>0$,
\begin{align}
\|\mu_t\|_{TV} &\leq \|\mu_0\|_{TV} e^{\alpha (2\MT{L})^{\gamma} (1+N) t} \label{pgcd2}
\\
\dst\int_{\R^+} xd\mu_t(x) &=\dst\int_{\R^+} xd\mu_0(x) \label{pgcd3}
\end{align}
where $N$ is defined in~\eqref{hyp1} and $L$ is defined in~\eqref{hyp2}. In particular such a solution is unique.
If moreover $\mu_t$ is non-negative (i.e. $\mu_t \in \clC(\R^+,\M^+(\R^+))$), then
\begin{equation}\label{pgcd11}
\supp(\mu_t) \subset [0,\MT{L}]. 
\end{equation}
\end{theo}

\begin{proof}
Consider $\mu_t \in \clC(\R^+,\M^+(\R^+))$ a non negative measure-valued solution to~\eqref{eq:frag} in the sense of Definition~\eqref{def:measure_valued_sol}. We start proving property~\eqref{pgcd11}. To this end we first notice that
\begin{align*}
\alpha x^\gamma \mu _t(x)=\MT{\alpha}\mu _t(x)\int _0^x\frac {y} {x}x^{\gamma -1}\kappa \left(\frac {y} {x}\right)dy.
\end{align*}
Then in the right-hand side of~\eqref{ppcm1} we write
\begin{multline*}
\int_{\R^+} d\mu_s(x) \alpha x^{\gamma}\left(- \varphi(x) + \dst\int_{0}^{1} d\kappa(z)\varphi(xz)\right)\\
=\int _0^\infty\int _0^x b(x, y)\, y\left(\frac {\varphi (y)} {y}- \frac {\varphi (x)} {x}\right)dyd\mu _s(x).
\end{multline*}
where
\begin{align}\label{bxy}
b(x,y)=\MT{\alpha x^{\gamma-1}\kappa \left(
\frac{y}{x}\right)}.
\end{align}
Consider then the test function
\begin{align*}
\varphi (x)=
\begin{cases}
0 &\forall x\in [0, L]\\
x(x-L)\quad &\forall x\in [L, L+1]\\
x &\forall x\ge L+1.
\end{cases}
\end{align*}
Since $x\to \varphi (x)/(1+x)$ is bounded and non decreasing on $[0, \infty)$, by~\eqref{ppcm1}
\begin{align*}
\int_{\R^+} \varphi(x) d\mu_t(x)&= \int_{\R^+}
\varphi(x) d\mu_0(x) + \int_0^t \int _0^\infty\int _0^x b(x, y)\, y\left(\frac {\varphi (y)} {y}- \frac {\varphi (x)} {x}\right)dyd\mu _s(x) ds\\
& \le \int_{\R^+} \varphi(x) d\mu_0(x) =0,
\end{align*}
where the last inequality is justified since $\mu_t\geq 0$ for $t\geq 0$ and since $x\to \varphi(x)/x$ is non decreasing as well. Since $\varphi\ge 0$ and $ \mu _t\ge 0$ for all $t\ge 0$ it follows that for $\varphi (x)d\mu _t(x)=0$ for all $t>0$ and almost every $x>0$. Since by construction $\varphi (x)>0$ for all $x>\MT{L}$ we deduce that for every $t>0$, $\supp (\mu _t)\subset [0, \MT{L}]$.
Consider now $\mu_t \in \clC(\R^+,\M(\R^+))$ a measure-valued solution to~\eqref{eq:frag} in the sense of Definition~\eqref{def:measure_valued_sol} (not necessarily non negative). To prove the BV estimate~\eqref{pgcd2}, we use definition~\href{https://ahl.centre-mersenne.org/item/AHL_2024__7__621_0/}{(1.6)}, and take $\varphi \in \clC(\R^+)$ such that $\|\varphi\|_{\infty} \leq 1$. By~\eqref{ppcm1},
\begin{multline*}
\int_{\R^+} \varphi(x) d\mu_t(x)\\
=\MT{\int_{\R^+} \varphi(x) d\mu_0(x)+}\dst\int_0^t ds \dst\int_{\R^+} \alpha x^{\gamma} d\mu_s(x) \left(- \varphi(x) + \dst\int_{z=0}^{1} \varphi(xz)d\kappa(z)\right).
\end{multline*}
Let $\chi\in C([0, \infty))$ such that $||\chi||_\infty=1$, $\chi(x)=1$ for all $x\in [0, \MT{L}]$ and $\chi(x)=0$ for $x>2\MT{L}$ and consider the function defined as
\begin{equation*}
\psi(x) := \alpha x^{\gamma}\chi(x) \left(- \varphi(x) + \dst\int_{z=0}^{1} \varphi(xz)d\kappa(z)\right),
\end{equation*}
It satisfies $ \psi \in \clC(\R^+)$ and, since $\|\varphi\|_{\infty} \leq 1$, and $\supp(\chi)\subset [0, 2 \MT{L}]$,
\begin{equation*}
\sup_{0\,\le\,x\,\le\,L} |\psi\MT{(x)}|\leq \alpha (2 \MT{L})^{\gamma} (1+N).
\end{equation*}
Therefore,
\begin{equation*}
\int_{\R^+} \varphi(x) d\mu_t(x)\leq \|\mu_0\|_{TV}+ \alpha (2\MT{L})^{\gamma} (1+N) \dst\int_0^t \|\mu_s\|_{TV} ds,
\end{equation*}
which implies
\begin{equation*}
\|\mu_t\|_{TV}\leq \|\mu_0\|_{TV}+ \alpha (2\MT{L})^{\gamma} (1+N) \dst\int_0^t \|\mu_s\|_{TV} ds,
\end{equation*}
and Gronwall Lemma yields~\eqref{pgcd2}. Finally, mass conservation property~\eqref{pgcd3} is obtained by choosing $\varphi(x)=x$ in definition~\ref{def:measure_valued_sol}
and using the last statement of~\eqref{hyp1}.
\end{proof}

The following proposition provides us with a solution to the fragmentation equation when initial condition is a Dirac delta localized at $x=\ell$ in terms of a fundamental solution.

\begin{prop}[Fundamental solution rescaled]\label{prop:rescaling}
Fix $\ell>0$.
Assume that $\mu_t^{F}$ is a fundamental solution to~\eqref{eq:frag}, i.e. a solution to~\eqref{eq:frag} when the initial data is $\mu_0=\delta(x-1)$. Then, $\mu_t^{\ell} = T_{\ell} \# \mu_{\ell^{\gamma}t}^{F}$, with $T_{\ell}(x)= \ell x$, is a solution to~\eqref{eq:frag} with $\mu_0^{\ell}=\delta(x-\ell)$.
\end{prop}

\begin{proof}
We set $\mu_t^{\ell}:= T_{\ell} \# \mu_{\ell^{\gamma}t}^{F}$. Let us prove that $\mu_t^{\ell}$ is a solution to~\eqref{eq:frag} with initial condition $\mu_0^{\ell}=\delta(x-\ell)$ and conclude by uniqueness of the solution. First, $\mu_0^{\ell} = T_{\ell}\# \mu_0 = T_{\ell}\# \delta(x-1) =\delta(x-\ell)$. Then, we obtain that for all $\varphi\in \clC_c(\R^+)$
\begin{align*}
\dst\int_{\R^+} \varphi(x) d\mu_t^{\ell}(x) &= \dst\int_{\R^+} \varphi(x) d\left(T_{\ell}\# \mu_{\ell^{\gamma}t}^{F}\right)(x) \\
&= \dst\int_{\R^+} (\varphi\circ T_{\ell})(x) d \mu_{\ell^{\gamma}t}^{F}(x) = \dst\int_{\R^+} \varphi(\ell x) d \mu_{\ell^{\gamma}t}^{F}(x).
\end{align*}
Since $\mu_t$ is a measure-valued solution to~\eqref{eq:frag}, we have
\begin{multline*}
\dst\int_{\R^+} \varphi(\ell x) d \mu_{\ell^{\gamma}t}^{F}(x)
=
\int_{\R^+} \varphi(\ell x) d\mu_0(x)\\
+\a \dst\int_0^{\ell^{\gamma t}} \dst\int_{\R^+}
\left(-x^{\gamma} \varphi(\ell x)d\mu_s^{F}(x)+ \varphi(\ell x)\dst\int_x^{\infty}y^{\gamma-1} d\kappa\left(\dfrac{x}{y}\right)d\mu_s^{F}(y)\right)ds.
\end{multline*}
Let us treat each of the three terms of the sum above separately. The first term is
\begin{equation*}
\int_{\R^+} \varphi(\ell x) d\mu_0(x) =
\int_{\R^+} \varphi(x) d\mu_0^{\ell}(x).
\end{equation*}
The second term is treated using the change of variables $s=\ell^{\gamma}u$
\begin{equation*}
-\a \dst\int_0^{\ell^{\gamma} t} \dst\int_{\R^+} x^{\gamma} \varphi(\ell x)d\mu_s^{F}(x)ds= -\a \dst\int_0^{t} \dst\int_{\R^+}
\left(x\ell\right)^{\gamma} \varphi(\ell x)d\mu_{\ell^{\gamma}u}^{F}(x)du,
\end{equation*}
and then the change of variables $z=T_{\ell}(x)$ i.e. $d\mu_{\ell^{\gamma}u}^{F}(x)= d(T_{\ell} \#\mu_{\ell^{\gamma}u}^{F})(z)$
\begin{align*}
-\a \dst\int_0^{t} \dst\int_{\R^+}
\left(x\ell\right)^{\gamma} \varphi(\ell x)d\mu_{\ell^{\gamma}u}^{F}(x)du &= -\a \dst\int_0^{t} \dst\int_{\R^+} z^{\gamma} \varphi(z)d\left(T_{\ell} \# \mu_{\ell^{\gamma}u}^{F}\right)(z)du\\
&= -\a \dst\int_0^{t} \dst\int_{\R^+} z^{\gamma} \varphi(z)d\mu_u^{\ell}(z)du.\\
\end{align*}
For the third term we also use the change of variables $s=\ell^{\gamma}u$ followed by the change of variables $z=T_{\ell}(x)$ and to finish the change of variable $w=T_\ell (y)$ i.e. $d\mu_{\ell^{\gamma}u}^{F}(y)= d (T_{\ell} \#\mu_{\ell^{\gamma}u}^{F})(w)=d\mu_u^\ell(w)$ to get
\begin{multline*}
\a \dst\int_0^{\ell^{\gamma t}} \dst\int_{\R^+}
\varphi(\ell x)\dst\int_x^{\infty}y^{\gamma-1} d\kappa\left(\dfrac{x}{y}\right)d\mu_s^{F}(y)ds \\
=\a \dst\int_0^{t} \dst\int_{\R^+}
\varphi(z)\dst\int_z^{\infty}w^{\gamma-1} d\kappa\left(\dfrac{z}{w}\right)d\mu_u^{\ell}(w)du.
\end{multline*}
To summarize,
\begin{multline*}
\dst\int_{\R^+} \varphi(x) d\mu_t^{\ell}(x) =
\dst\int_{\R^+} \varphi(x) d\mu_0^{\ell}(x)\\
-\a \dst\int_0^{t} \dst\int_{\R^+}
\left(z^{\gamma} \varphi(z)d\mu_u^{\ell}(z)+
\varphi(z)\dst\int_z^{\infty}w^{\gamma-1} d\kappa\left(\dfrac{z}{w}\right)d\mu_u^{\ell}(w) \right)du.
\end{multline*}
Finally, since $t\to \mu_t^{F}$ is narrowly continuous, then $t\to \mu_t^{\ell}$ is narrowly continuous as well. This ends the proof of Proposition~\ref{prop:rescaling}.
\end{proof}

\begin{theo}[Existence of a solution to~\eqref{eq:frag} represented as a power series]\label{thm:power_serie}
For any fragmentation kernel $\kappa $ satisfying~\eqref{hyp1}, $\gamma\geq 0$ and $\mu_0$ satisfying~\eqref{hyp2}, there exists a measure-valued solution $\mu_t\in\clC(\R^+,\M(\R^+))$ to~\eqref{eq:frag} in the sense of Definition~\ref{def:measure_valued_sol}. This solution is given by the following everywhere convergent series
\begin{equation}\label{representation_solution}
\mu_t = e^{-\a x^{\gamma}t} \mu_0 + \sum_{n=0}^{\infty} (\a t)^n
\dst\int_0^{\infty} \ell^{n\gamma}a_n\left(\dfrac{x}{\ell}\right)\mu_0(\ell) \dfrac{d\ell}\ell{},
\end{equation}
where the sequence $a_n$ is defined as follows for $x\in [0,1]$,
\begin{equation}\label{induction}
\begin{split}
a_0(x) &=0, \\
 a_{n+1}(x) &= \dfrac{1}{n+1}\left(-x^{\gamma}a_n(x)+ \dst\int_x^{\infty} y^{\gamma-1} \kappa \left(\dfrac{x}{y} \right)a_n(y)dy + \kappa(x) \dfrac{(-1)^n}{n!}\right).
\end{split}
\end{equation}
In particular,
\begin{equation}\label{weakform}
\dst\int_0^{\infty} \varphi(x) d\mu_t(x) = \dst\int_0^{L} \dst\int_0^{1} \varphi(\ell x) d \mu_{t\ell^{\gamma}}^{F} (x)d\mu_0(\ell),\qquad \forall \varphi\in \clC(\R^+),
\end{equation}
and
\begin{equation}
\supp(\mu_t)\subset \supp(\mu_0), \qquad \text{for all } t>0.
\end{equation}

\end{theo}



\begin{rema}\label{remark1}
We emphasize that in~\eqref{induction} and~\eqref{representation_solution}, $\kappa$, $\mu_0$ and $a_n$ may be measures. In this case, the Mellin convolution product has to be understood in the sense of measures. For instance, the formula~\eqref{representation_solution} means
\begin{align}
\mu_t &= e^{-\a x^{\gamma}t} \mu_0 + \sum_{n=0}^{\infty} (\a t)^n a_n \ast b_n \label{remark1E1}
\\
 db_n(y) &=y^{n\gamma } d\mu _0(y), \label{remark1E2}
\end{align}
where the Mellin convolution product $\MT{\ast}$ is defined in Definition~\href{https://ahl.centre-mersenne.org/item/AHL_2024__7__621_0/}{1.2}.
\end{rema}

\begin{proof}[Proof of Theorem~\ref{thm:power_serie}]\leavevmode
\begin{step}{{A power series representation for the fundamental solution}}\label{prftheo2.4.stp1}
In this step, we prove that
\begin{align}\label{shape_u}
&\mu_t^{F}(x)=e^{-\a t}\delta(x-1) +v_t, 
\end{align}
where
\begin{align}\label{serie}
v_t= \sum_{n=1}^{\infty} (\a t)^n a_n \in \clC((0,T),\M(\R^+)), \qquad v_t(0)=0, \qquad v_t\perp \delta (x-1),
\end{align}
is a fundamental solution to~\eqref{eq:frag}, and we prove that for all $t>0$, this solution satisfies
\begin{align}
& \supp\left(\mu^F_t\right)\subset[0, 1].
\end{align}
Fix $T>0$. Assume that $\mu_t^F$ is a fundamental solution to~\eqref{eq:frag} on $[0,T]\times \R^+$. We recall that $\mu \perp \nu$ if there exists $E \in \clB(\R)$ such that $\mu(\R)= \mu(E)$ and $ \nu(E)=0$. The Radon--Nikodym decomposition guarantees that for all $t>0$, $\mu_t^F$ can be decomposed as
\begin{equation}\label{shape_u1}
\mu_t^{F} = A(t) \delta (x-1) +v_t,
\end{equation}
where $v_t \perp \delta(x-1)$, $A(0)=1$ and $v_0(x)=0$. We plug~\eqref{shape_u1} into~\eqref{eq:frag} and get
\begin{multline*}
A'(t) \delta(x-1)+ \p_tv_t(x)\\
=- A(t) \a x^{\gamma}\delta(x-1) - v_t(x) \alpha x^{\gamma} +
\alpha \dst\int_x^{\infty} y^{\gamma-1} \kappa\left(\dfrac{x}{y} \right)
\left(A(t) \delta(y-1) +dv_t(y)\right)
\end{multline*}
which is
\begin{multline*}
A'(t) \delta(x-1)+ \p_tv_t(x) \\
=- \a A(t) \delta(x-1) - v_t(x) \alpha x^{\gamma} +
\alpha \kappa(x) A(t)+ \alpha \dst\int_x^{\infty} y^{\gamma-1} \kappa\left(\dfrac{x}{y} \right) dv_t(y).
\end{multline*}
By identification, we get that necessarily
\begin{equation}\label{eq:Aandv}
\left\{
\begin{aligned}
A'(t)&= - \alpha A(t), \quad A(0)=1, \\
\p_t v_t(x)&= - \alpha x^{\gamma}v_t(x)+\alpha \dst\int_x^{\infty} y^{\gamma-1} \kappa\left(\dfrac{x}{y}\right)dv_t(y) + \alpha \kappa(x) A(t), \quad v_0(x)=0.
\end{aligned}
\right.
\end{equation}
The first line gives $ A(t)= e^{-\a t}$. This proves that $\mu_t^F$ is necessarily equal to $e^{-\alpha t} +v_t$, where $v_t$ satisfies the second line of~\eqref{eq:Aandv}. Now let us verify that the series~\eqref{serie} converges in $\clC((0,T),\M(\R^+))$. Since $(\M(\R^+),\|.\|_{TV})$ is a Banach space, it is enough to prove the normal convergence of the series~\eqref{serie}. We first claim that
\begin{equation}\label{ineq_serie}
\| a_{n+1} \|_{TV} \leq \frac{1}{n+1} \left((N+1)\| a_{n} \|_{TV}+ \f{N}{n!} \right),
\end{equation}
This comes directly from the induction formula~\eqref{induction} since $x\in [0,1]$ implies
\begin{equation*}
\| x^{\gamma} a_{n} \|_{TV} \leq \| a_{n} \|_{TV},
\end{equation*}
and
\begin{align*}
\left\| \dst\int_x^{\infty} y^{\gamma-1} \kappa \left(\frac{x}{y} \right)a_n(y)dy\right\|_{TV} &= \dst\int_0^{\infty} \left| \dst\int_x^{\infty} y^{\gamma-1} \kappa \left(\frac{x}{y} \right)a_n(y)dy \right| dX\\
&\leq \dst\int_0^{\infty}\dst\int_x^{\infty} \left| y^{\gamma-1} \kappa \left(\frac{x}{y} \right)a_n(y)\right|dy dx\\
&=\dst\int_0^{1}\dst\int_0^{\infty} \left| y^{\gamma} \kappa \left(z \right)a_n(y)\right|dy dz\\
&\leq \dst\int_0^{1} \left|\kappa \left(dz \right)\right| \dst\int_0^{\infty} y^\gamma\left|a_n(y)\right|dy
\leq N \| a_{n} \|_{TV},
\end{align*}
and finally
\begin{equation*}
\left\| \kappa(x) \frac{(-1)^n}{n!}\right\|_{TV} = \f{N}{n!}.
\end{equation*}
We deduce from~\eqref{ineq_serie} that for all $n\in\N,$
\begin{equation}\label{u_n}
\| a_n \|_{TV} \leq \f{(N+2)^n}{n!},
\end{equation}
hence the \Mag{normal} convergence of the series~\eqref{serie} in $(\M(\R^+),\|.\|_{TV})$ for all $t>0$. We prove~\eqref{u_n} by induction:~\eqref{u_n} is true for $n=0$ and $n=1,$ and if it is satisfied for $n\geq 1$ we have

\begin{align*}
\Vert a_{n+1}\Vert_{TV} &\leq \f{1}{n+1} \left((N+1)\Vert a_n\Vert_{TV} + \f{N}{n!}\right)\\
&\leq \f{1}{(n+1)!} \f{(N+1)(N+2)^{n} + N}{n!} \leq \f{(N+2)^{n+1}}{(n+1)!}.
\end{align*}
Then the series $v_t$ defined in~\eqref{serie} converges in the Banach space $\clC((0,T),\M(\R^+)).$ Since from the induction rule~\eqref{induction} $\supp(a_n)\subset [0, 1]$ for all $n\ge 0$, it follows that $\supp(v_t)\subset [0, 1]$. We prove then, using the differentiation rule of power series in a Banach space, that the power series~\eqref{serie} is a solution to the second line of~\eqref{eq:Aandv}. We have
\begin{equation*}
\dfrac{d}{dt} v(t,x) = \dfrac{d}{dt} \sum_{n=1}^\infty (\a t)^n a_n=
\a \sum_{n=1}^\infty n (\a t)^{n-1} a_n = \a \sum_{n=0}^\infty (n+1) (\a t)^n a_{n+1}.
\end{equation*}
Using the induction hypothesis~\eqref{induction} we get
\begin{multline*}
\frac{d}{dt} v(t,x)\\
\begin{aligned}
&=  \a \sum_{n=0}^\infty (\a t)^{n}\left(-x^{\gamma}a_n(x) + \dst\int_x^{\infty} y^{\gamma-1} \kappa \left(\dfrac{x}{y} \right)a_n(y)dy + \kappa(x) \dfrac{(-1)^n}{n!}\right)\\
& = - \a x^{\gamma} \sum_{n=0}^\infty (\a t)^{n} a_n +\a \dst\int_x^{\infty} y^{\gamma-1} \kappa \left(\dfrac{x}{y} \right)\left(\sum_{n=0}^\infty (\a t)^{n} a_n(y)\right) dy +\a \sum_{n=0}^\infty (-\a t)^{n}\kappa(x)\\
& = - \a x^{\gamma} v(t,x) +\a \dst\int_x^{\infty} y^{\gamma-1} \kappa \left(\dfrac{x}{y} \right)v(t,y)dy +\a e^{-\a t} \kappa(x),
\end{aligned}
\end{multline*}
which is~\eqref{eq:Aandv}. The property of the support $\supp(\mu _t^F)\subset [0, 1]$ follows from the hypothesis on the support of $\kappa$ and by inspection of formulas~\eqref{serie} and~\eqref{induction}.
\end{step}

\begin{step}{{A power series representation of a solution with a generic initial condition}}\label{prftheo2.4.stp2}
By the classical superposition principle, if
\begin{equation}\label{integralmu1}
\mu_t (x)= \int _0^\infty \mu _0(\ell)\mu _t^\ell\left(x \right)d\ell
\end{equation}
converges in $\clC((0,T),\M(\R^+))$ (where $\mu_t^{\ell} $ is the scaled fundamental solution obtained in Proposition~\ref{prop:rescaling} from $\mu_t^F$ the fundamental solution obtained in Step~\ref{prftheo2.4.stp1}), $\mu _t$ will be a solution to the fragmentation equation with initial condition $\mu_0 \in \M(\R^+)$. Notice that we have the following equality for the integral~\eqref{integralmu1}:
\begin{align}\label{seriemu1}
\int_0^{\infty} \mu_0(\ell) \mu_t^{\ell} \left({x} \right) d\ell=e^{-\alpha x^\gamma t}\mu _0+\dst\int_0^{\infty}\sum_{n=0}^{+\infty} (\a t)^n
\ell^{n\gamma}a_n\left(\dfrac{x}{\ell}\right)\mu_0(\ell) \dfrac{d\ell}\ell{}.
\end{align}
Since for every $n$, by~\eqref{u_n}
\begin{align*}
\sum_{n=0}^{m} \left\| (\a t)^n
\dst\int_0^{\infty} \ell^{n\gamma}a_n\left(\dfrac{x}{\ell}\right)\mu_0(\ell) \dfrac{d\ell}\ell{}\right\| _{ TV } &\le \|\mu _0\| _{ TV }\sum_{n=0}^{m} (\a t)^n\|a_n\| _{ TV } L^{\gamma n} \\
&\le \|\mu _0\| _{ TV } \sum_{n=0}^{m} (\a t)^n \f{(N+2)^n}{n!}L^{\gamma n},
\end{align*}
the series in the right-hand side of~\eqref{seriemu1} converges absolutely in the Banach space $\clC((0,T),\M(\R^+))$ for all $T>0$. The integral~\eqref{integralmu1} is then absolutely convergent and defines a solution to the fragmentation equation with initial condition $\mu_0 \in \M(\R^+)$ and for $t\in[0,T]$. Property~\eqref{weakform} follows then from the definition of $\mu^F _{ t\ell^\gamma}$. Since $\supp(a_n)\subset [0, 1]$ for every $n\ge 0$, it follows from~\eqref{representation_solution} that $\supp(\mu _t)\subset [0, L]$ for all $t>0$. Using a classical diagonal argument, and since the property on the support of the solution does not depend on $T$, the power series defines a solution in $\clC({\R}^+,\M(\R^+))$ This ends the proof of Theorem~\ref{thm:power_serie}].\qedhere
\end{step}\let\qed\relax
\end{proof}



\begin{theo}[Non negativity of the power series solution]\label{thm:positivity}
Assume the fragmentation kernel $\kappa$ satisfies~\eqref{hyp1} with $\gamma \geq 0$ and take $\mu_0$ that satisfies~\eqref{hyp2}. Then, the power series solution~\eqref{representation_solution} to the fragmentation equation~\eqref{eq:frag} is non-negative.
\end{theo}

\begin{rema}
Up to our knowledge, no proof of positivity for the fragmentation equation is available in the literature in the case where either the initial condition $\mu_0$ is a measure or the fragmentation kernel $\kappa$ is a measure. In our case, both are measures.
\end{rema}

\begin{proof}
We first prove the non negativity of the fundamental solution $\mu_t^F$ defined by~\eqref{shape_u}~\eqref{serie} using an approximation argument. Consider to this end the function $\chi$ such that:
\begin{align*}
&\chi\in \clC(0, \infty),\;\chi\ge 0,\;\|\chi\|_\infty=1, \quad \chi(x)=1\quad\forall x\in [0, 2L];\quad
\chi(x)=0\;\forall x>3L,
\end{align*}
and a mollifier $\theta\in \clC(0, \infty)$, such that $\theta \ge 0$, $\supp \theta\subset [0, 1]$ and the sequence $\theta _m(x)=m\theta (m(x-1))$. Let us then denote by $\kappa_m$ a regularization of the fragmentation kernel
\begin{equation*}
\kappa _m=\kappa \ast \theta_m,
\end{equation*}
by $a$ a regularization of the fragmentation rate
\begin{equation*}
a(x)=x^{\gamma}\chi(x),
\end{equation*}
and by $\theta_k$ a regularization of the initial condition $\delta(x-1)$. By construction $\|a\|_\infty. \linebreak\le (3L)^\gamma$ and for each $m\ge 1$, $\kappa _m$ is a regular function satisfying
\begin{equation}
\label{kmBL2}
\begin{split}
\supp[\kappa_m]\subset [0, 2], \quad \lim_{m\,\to \,\infty}\|\kappa _m-\kappa\|_{BL} &=0, \\
\|\kappa _m\|_{TV}\le \|\MT{\kappa}\|_{TV}\|\theta_m\|_{TV} &\le N.
\end{split}
\end{equation}
Consider for every $m\ge 1$ the sequence of functions $\{a_{n, m}\}_{n\,\in\,\N,\,m\,\in\,\N}$ defined as follows,
\begin{multline}\label{inductionB}
 a_{0, m}(x) =0,\\
a_{n+1, m}(x) = \dfrac{1}{n+1}\left(-a(x)a_{n, m}(x)+ \dst\int_{\MT{0}}^{\infty} a(y) \kappa_m \left(\dfrac{x}{y} \right)a_{n, m}(y)\MT{\dfrac{dy}{y}} + \kappa_m(x) \dfrac{(-1)^n}{n!}\right).
\end{multline}
It immediately follows, for all $n\ge 1, m\ge 1$,
\begin{align*}
a_{n, m}\in \clC([0, \infty)),\,\,\supp[a_{n, m}]\subset[0, 2]
\end{align*}
and then, the sequence $\{a_{n, m}\}_{n\,\in\,\N,\,m\in\,\N}$ satisfies also,
\begin{multline}\label{inductionC}
a_{n+1, m}(x)\\
= \frac{1}{n+1}\left(-x^{\gamma}a_{n, m}(x)+ \dst\int_{\MT{0}}^{\infty} y^{\gamma} \kappa_m \left(\dfrac{x}{y} \right)a_{n, m}(y)\frac{dy}{y} + \kappa_m(x) \dfrac{(-1)^n}{n!}\right).
\end{multline}
Since $\supp[a_{n, m}]\subset [0, 2]$ it follows that
\begin{equation*}
\left\|x\to x^\gamma a_{n, m}(x)\right\|_{TV}\le 2^{\gamma}\left\|a_{n, m}\right\|_{TV}.
\end{equation*}
Then, for every $m\ge 1$ fixed, as for the proof of~\eqref{u_n} it follows now that
\begin{equation}\label{u_nm}
\| a_{n, m} \|_{TV} \leq \f{2^{\gamma n}(N+2)^n}{n!},\,\,\forall n\ge 1, \, \forall m\ge 1.
\end{equation}
\setcounter{stepcount}{0}
\begin{step}{{The solution $u_{k,m}$ to the regularized problem is non negative.}}\label{prftheo2.6.stp1}
Consider the regularized problem ($m<\infty$ and $k<\infty$)
\begin{equation}\label{pgcd1}
\left\{
\begin{aligned}
&\dfrac{\p}{\p t} u_{k, m}(t, x) =-\a a(x) u_{k, m}(t, x) + \alpha \dst\int_{x}^{\infty}\kappa_m\left(\dfrac{x}{y} \right)a(y) u_{k, m}(t, y)\frac{dy}{y},\\
&\MT{u_{k,m}(0,x)=\theta_k(x)}.
\end{aligned}
\right.
\end{equation}
We define for $k\geq 1$ and $m\geq 1$ the sequence of functions
\begin{align}\label{pgcd8C}
u_{k, m}(t, x) = e^{-\a x^{\gamma}t} \theta_k +
\sum_{n=0}^{\infty} (\a t)^n
\dst\int_0^{\infty} \ell^{n\gamma}a_{n, m}\left(\dfrac{x}{\ell}\right) \theta_k(\ell) \dfrac{d\ell}\ell{},
\end{align}
For $k\geq 1$ and $m\geq 1$, the series is absolutely convergent. Moreover, by construction $\supp[u_{k, m}]\subset [0, 2]$. Then, $u_{k, m}$ satisfies~\eqref{pgcd1}. By~\cite[Lemma~3]{Melzak}, of which the equation~\eqref{pgcd1} and the initial data $\theta_k$ satisfy the hypothesis, the Cauchy problem for~\eqref{pgcd1} with initial data $\theta_k$ possesses a global solution bounded, continuous, non negative, analytic in $t$ for each $x>0$ and integrable in $x$ for every $t>0$. Moreover, by construction, for all $T>0$ and $t\in (0, T)$,
\begin{equation*}
\begin{aligned}
\dst\int\left|u_{km}(t,x)\right| dx& =
\dst\int_0^{\infty} e^{-\alpha x^{\gamma}t}\theta_k(x)dx+ \dst\int_0^{\infty}\sum_{n=0}^{\infty} (\a t)^n
\dst\int_0^{\infty} \ell^{n\gamma}\left|a_{n, m}\left(\dfrac{x}{\ell}\right)\right| \theta_k(\ell) \dfrac{d\ell}\ell dx\\
&\leq 1+ \sum_{n=0}^{\infty} (\a t)^n 2^{n\gamma} \dfrac{(N+2)^n}{n!}<\infty.
\end{aligned}
\end{equation*}
Therefore, by~\cite[Theorem~1]{Melzak}, the function $u_{k,m}$ is the unique solution of \eqref{pgcd1}
that is integrable with respect to $x$.
Moreover, this solution is non negative.
\end{step}


\begin{step}{{limit $k\to \infty$.}}\label{prftheo2.6.stp2}
Consider the problem ($m<\infty$)
\begin{equation}\label{eq:fragReg}
\begin{aligned}
\dfrac{\p}{\p t} u_{m}(t, x) &=-\a a(x) u_{m}(t, x) + \alpha \dst\int_{x}^{\infty}\kappa_m\left(\dfrac{x}{y} \right)a(y) u_{m}(t, y)\frac{dy}{y},\\
u_{m}(0, x) &= \delta (x-1).\
\end{aligned}
\end{equation}
We define the sequence of measures $u_m$ as
\begin{align}\label{pgcd8}
u_{m}(t, x) = e^{-\a t} \delta (x-1) +
\sum_{n=0}^{\infty} (\a t)^n a_{n, m}
\end{align}
The series in~\eqref{pgcd8} is absolutely convergent in $TV$ norm for every $m\ge 1$ and it defines a measure $u_m\in \M(\R_+)$ such that
\begin{align}
\|u_m(t)\|_{TV}\le \exp\left(\alpha t 2^{\gamma }(N+2)\right)\MT{+e^{-\alpha t}}.
\end{align}


The measure $u_{m}(t,.)$ satisfies~\eqref{eq:fragReg} but since $\supp[u_m(t)]\subset [0, 2]$ it also satisfies,
\begin{align*}
&\dfrac{\p}{\p t} u_{m}(t, x) =-\a x^{\gamma} u_{m}(t, x) + \alpha \dst\int_{x}^{\infty}\kappa_m\left(\dfrac{x}{y} \right) y^{\gamma-1} u_{m}(t, y)dy.
\end{align*}
We claim now that $(u_{km})_{k\,\geq\,0}$ converges weakly towards $u_m$. Indeed, on one hand, for all $\varphi \in \clC_c(\R^+)$,
\begin{equation}\label{fleur1}
\lim\limits_{k\,\to\,\infty} \dst\int_0^{\infty} e^{-\alpha x^{\gamma} t} \theta_k(x) \varphi(x)dx = e^{-\alpha t} \varphi(1),
\end{equation}
and on the other hand, let us notice that for each $n\ge 1$ and $m\ge 1$ fixed,
\begin{equation*}
\lim_{k\,\to\,\infty} \int_0^{\infty} \ell^{n\gamma}a_{n, m}\left(\dfrac{x}{\ell}\right) \theta_k(\ell) \dfrac{d\ell}\ell{}=a_{n, m}(x), \quad x\in [0,2].
\end{equation*}
Moreover, for all $k\ge 1$, $n\ge 1$ and $m\ge 1$
\begin{align*}
\int_0^{\infty} \ell^{n\gamma}a_{n, m}\left(\dfrac{x}{\ell}\right) \theta_k(\ell) \dfrac{d\ell}\ell{}=
\int _0^\infty \left(\frac{x}{y}\right)^{n\gamma-1}\!\!\!a_{n, m}(y)k\, \theta\left(k\frac{x}{y}\right) dy.
\end{align*}
Therefore, if
\begin{align*}
\psi _k (z)= z^{n\gamma-1}k \theta (k z)
\end{align*}
then $\psi _k\in \clC([0, \infty)$ and $\psi _k\ge 0$. Since $\supp \theta \subset [0, 1]$ and $\|\theta\|_\infty\le 1$, for all $n$ such that $n\gamma >2$,
\begin{align*}
\sup_{z\,>\,0}\psi _k(z)=\sup_{z>0 }z^{n\gamma-1}k\theta(k z)=
\sup_{z\,\in\,\left[0, k^{-1}\right]}z^{n\gamma-2}(kz)\theta(k z)\le 1
\end{align*}
It follows that for all $m\ge 1$, $n\ge 1$, $k\ge 1$ and $x>0$,
\begin{align*}
\left|\int_0^{\infty} \ell^{n\gamma}a_{n, m}\left(\frac{x}{\ell}\right) \theta_k(\ell) \frac{d\ell}\ell{}\right|
\le \|a_{n, m}\|_{TV}
\end{align*}
It follows by the Lebesgue's convergence that
\begin{equation}\label{fleur2}
\lim\limits_{k\,\to\,\infty} \sum_{n=0}^{\infty} (\a t)^n
\dst\int_0^{\infty} \ell^{n\gamma}a_{n, m}\left(\dfrac{x}{\ell}\right) \theta_k(\ell) \dfrac{d\ell}\ell{} = \sum_{n=0}^{\infty} (\a t)^na_{n,m}(x), \quad x\in[0,2],
\end{equation}
and then, combining~\eqref{fleur1} with~\eqref{fleur2} gives us
\begin{align*}
\lim_{k\,\to\,\infty} \int _0^\infty \varphi (x) u_{k, m}(t, x)dx =\int _0^\infty \varphi (x) u_m(t, x)dx,\quad \varphi\in \clC_c(\R^+),
\end{align*}
and then
\begin{equation}
u_m\ge 0.
\end{equation}
\end{step}
\begin{step}{{limit $m\to \infty$.}}\label{prftheo2.6.stp3}
We prove here that $u_m$ converges weakly towards $\mu_t^F$. To do so, we prove by induction that
\begin{equation}
\left\|a_{n,m}-a_n\right\|_{BL}\underset{m\,\to\,\infty}{\to} 0.
\end{equation}
For $n=1$, $a_{1, m}=\kappa _m$, $a_1=\kappa$, and by construction $\|\kappa _m-\kappa \|_{BL} \underset{m\,\to\,\infty}{ \to} $. Assume then $\|a_{n, m}-a_n\|_{BL}\to 0$. In order to prove that the same property holds for the sequence $\{a_{n+1, m}\}_{m\,\in\,\N}$ it is sufficient to prove
\begin{align}\label{convmeme}
\lim_{m\,\to\,\infty}\left\|\int_x^{\infty} y^{\gamma-1} \kappa_m \left(\dfrac{x}{y} \right)a_{n, m}(y)dy - \int_x^{\infty} y^{\gamma-1} \kappa \left(\dfrac{x}{y} \right)a_{n}(y)dy\right\|_{BL}= 0.
\end{align}
If, for the sake of notation we define the functions $\tilde a_{n, m}$ and $\tilde a_{n}$ as
\begin{align*}
\tilde a_{n, m}(\ell) &=\ell^{\gamma}a_{n, m}(\ell),\qquad
\tilde a_{n}(\ell) &=\ell^{\gamma}a_{n}(\ell),
\end{align*}
then property~\eqref{convmeme} reads,
\begin{align}\label{convmemeB}
\lim_{m\,\to\,\infty}\left\|\tilde a_{n, m}\ast \kappa _m- \tilde a_{n}\ast \kappa \right\|_{BL}= 0.
\end{align}
Notice indeed that, for all test function $\varphi $ such that $\|\varphi \|_\infty\le 1$ and $\|\varphi '\|_\infty\le 1$,
\begin{multline}\label{comvmeme2}
\int_0^\infty \varphi (x)
\big(\tilde a_{n, m}\ast \kappa _m(x)- \tilde a_{n}\ast \kappa(x)\big)dx=\\
\int_0^\infty \varphi (x)
\left (\tilde a_{n, m}- \tilde a_{n}\right)\ast \kappa _m(x)dx
+\int_0^\infty \varphi (x)
\left (\kappa _m- \kappa \right)\ast\tilde a_{n}(x)dx
\end{multline}
The two terms in the right-hand side of~\eqref{comvmeme2} may be bounded with the same arguments. Consider for example the first.
\begin{align*}
\left|\int_0^\infty \varphi (x)
\left(\tilde a_{n, m}- \tilde a_{n}\right)\ast \kappa _m(x)dx\right|=\left|\int _0^\infty \varphi (x)\int _0^\infty
\left (\tilde a_{n, m}- \tilde a_{n}\right)\left(\frac{x}{y}\right)\kappa _m(y)\frac{dy}{y}dx\right|\\
=\left|\int _0^2 \kappa _m(y) \int _0^\infty
\left (\tilde a_{n, m}- \tilde a_{n}\right)\left(\frac{x}{y}\right)\varphi (x)dx\frac{dy}{y}\right|\\
\le \|\kappa_m\|_{ TV }
\sup_{y\,\in\,[0, 2]}
\left|\int _0^\infty
\left (\tilde a_{n, m}- \tilde a_{n}\right)\left(z\right)\varphi (zy)dz \right|
\end{align*}
For each $y\in[0, 2]$,
\begin{align*}
\left|\int _0^\infty
\left (\tilde a_{n, m}- \tilde a_{n}\right)\left(z\right)\varphi (zy)dz \right|\le \|\tilde a_{n, m}- \tilde a_{n}\|_{BL}
\left(\|\varphi \|_\infty+2\|\varphi'\|_\infty \right)
\end{align*}
from where, by~\eqref{kmBL2}
\begin{align*}
\left\|
\left (\tilde a_{n, m}- \tilde a_{n}\right)\ast \kappa _m\right\|_{BL}
\le 3 N \|\tilde a_{n, m}- \tilde a_{n}\|_{BL}.
\end{align*}
A similar arguments shows, using~\eqref{u_n},
\begin{align*}
\left\|
\left (\kappa _m- \kappa \right)\ast \tilde a_n \right\|_{BL}\le 3 \|k_m-k\|_{BL }\|\tilde a_n\|_{TV}\le 3 \f{(N+2)^n}{n!}\|k_m-k\|_{BL }
\end{align*}
and then, \eqref{convmeme} holds true.

Now, for any $\varphi\in \clC^1([0, \infty))$, such that $\|\varphi \|_\infty+\|\varphi '\|_\infty <\infty$, by definition of the measure $u_m$
\begin{align*}
\int _0^\infty u_m(t, x)\varphi (x)dx&= e^{-\alpha t}\varphi (1)+\int _0^\infty
\left (\sum_{n=0}^{\infty} (\a t)^n a_{n, m}(x)\right) \varphi (x)dx\\
&=e^{-\alpha t}\varphi (1)+ \sum_{n=0}^{\infty} (\a t)^n \int _0^\infty a_{n, m}(x) \varphi (x)dx.
\end{align*}
Since for every $n\ge 1$,
\begin{equation*}
\lim _{m\,\to\,\infty} \int _0^\infty a_{n, m}(x) \varphi (x)dx=\int _0^\infty a_{n}(x) \varphi (x)dx
\end{equation*}
and
\begin{align*}
(\a t)^n \int _0^\infty \left|a_{n, m}(x) \varphi (x)\right|dx& \le (\alpha t)^n\|\varphi \|_\infty \|a_{n, m}\|_{TV}\\
&\le (\alpha t)^n\|\varphi \|_\infty\f{2^{\gamma n}(N+2)^n}{n!}
\end{align*}
one has,
\begin{align*}
\lim _{m\,\to\,\infty} \int _0^\infty\varphi (x)u_{m}(t, x)dx&= e^{-\alpha t}\varphi (1)+ \sum_{n=0}^{\infty} (\a t)^n \int _0^\infty a_{n}(x) \varphi (x)dx=\int _0^\infty\varphi (x)d\mu _t^F
\end{align*}
and it follows
\begin{equation*}
\mu _t^F\ge 0.
\end{equation*}
This ends the proof of Theorem~\ref{thm:positivity}.\qedhere
\end{step}\let\qed\relax
\end{proof}

\begin{thebibliography}{99}
\bibitem[Mel57]{Melzak} Z. A. Melzak, ``A Scalar Transport Equation'', \emph{Trans. Am. Math. Soc.} \textbf{85} (1957), no. 2, 547--560.
\end{thebibliography}
\end{document}
